A central design goal is that scientific operations are executed through explicit computational tools, while the language model focuses on translating intent into a validated sequence of actions.
This aligns with research on tool-augmented agents, where external actions provide verifiable intermediate results and reduce reliance on unconstrained text generation \citep{Yao2022ReAct,Karpas2022MRKL}.

\subsection{Capability categories}
The tool layer operationalizes common workflows in ocean and environmental data analysis:
\begin{itemize}[leftmargin=*]
  \item \textbf{Discovery and resolution:} identify candidate datasets and variables, interpret metadata, and disambiguate naming.
  \item \textbf{Data access:} retrieve spatiotemporal subsets with user-defined bounds and sampling options.
  \item \textbf{Analysis:} compute derived products such as climatologies and summary statistics.
  \item \textbf{Integration:} colocalize datasets by joining observations and gridded products within spatial/temporal tolerances.
  \item \textbf{Visualization:} generate quick-look figures (maps, time series) that summarize retrieved subsets.
  \item \textbf{Reproducible code:} optionally emit executable code that reproduces the same operations in the CMAP SDK.
\end{itemize}

\subsection{Representative tool outputs}
Many interactions produce both a narrative response and scientific artifacts (data files, figures).
Table~\ref{tab:capabilities} summarizes the primary capability classes and expected outputs.

\begin{table}[h]
\centering
\small
\begin{tabular}{@{}p{0.34\linewidth}p{0.61\linewidth}@{}}
\toprule
\textbf{Capability} & \textbf{Typical outputs} \\
\midrule
Catalog discovery and variable resolution & Candidate datasets/variables with brief justification grounded in catalog context \\
Spatiotemporal subsetting & Downloadable data subset (e.g., Parquet/CSV) with provenance notes \\
Quick-look visualization & Map plots or time-series figures derived from retrieved subsets \\
Climatology and summaries & Climatology fields or seasonal cycles; summary tables/figures \\
Colocalization / dataset integration & Joined tables for multi-dataset analysis; tolerance settings and diagnostics \\
Code emission (optional) & Executable SDK code reproducing the performed operations \\
\bottomrule
\end{tabular}
\caption{Capability-oriented summary of \cmapagent. The underlying implementation exposes these functions through structured tool contracts; the main text emphasizes scientific functionality rather than specific software frameworks.}
\label{tab:capabilities}
\end{table}

\subsection{Iterative planning and tool execution loop}
At the core of \cmapagent\ is a per-turn runner that alternates between (i) requesting a structured plan from the language model and (ii) executing the requested tool calls.
The runner enforces a strict JSON protocol with two message types: \texttt{tool\_call} (a list of tool invocations) and \texttt{final} (the user-facing response).
It also implements practical guardrails that matter for scientific use: a configurable tool-call budget, caching of identical tool calls to avoid accidental loops, sanitization of co-localization tolerances when they were not specified by the user, and aggregation of tool-produced artifacts (e.g., plots and data extracts) into the final response.
Algorithm~\ref{alg:execute_plan} summarizes this control flow in compact pseudocode format.

\begin{algorithm}[!ht]
\caption{Per-turn agent runner (conceptual) implementing an iterative plan--execute--observe loop with key guardrails: strict structured outputs, bounded tool usage, per-turn deduplication, co-localization tolerance sanitization, and aggregation of tool-produced artifacts and optional reproducible code into the final response.}
\label{alg:execute_plan}
\small
\SetAlgoNlRelativeSize{-1}
\DontPrintSemicolon
\KwRequire{system prompt; conversation history; user message; execution context; tool-call budget $B$}
\KwResult{final response (message + optional code + artifact references), and tool trace for provenance}

$M \leftarrow [\textsf{SYSTEM}(\texttt{system\_prompt})] + \textsf{conversation} + [\textsf{USER}(\texttt{user\_message})]$\;
$T \leftarrow [\ ]$; $A \leftarrow [\ ]$; $C \leftarrow [\ ]$; $\mathcal{H} \leftarrow \emptyset$ \tcp*[r]{trace, artifacts, code, dedup cache}\;
$b \leftarrow \max(0,B)$; $r_{\textsf{json}}\leftarrow 0$; $r_{\textsf{force}}\leftarrow 0$\;
\If{$b = 0$}{Append to $M$ an instruction that tools are disabled and the planner must return \texttt{final}.}\;

\Repeat{a \texttt{final} response is produced}{
  $\texttt{resp} \leftarrow \textsc{LLMComplete}(M)$\;
  \If{\textsc{StrictJSONWithType}(\texttt{resp}) is false}{
    \If{$r_{\textsf{json}}<2$}{
      $r_{\textsf{json}}\leftarrow r_{\textsf{json}}+1$; Append a JSON-only repair instruction to $M$; \textbf{continue}\;
    }
    \KwRet{\textsc{FinalizeFromText}(\texttt{resp},A,C), $T$}\;
  }

  $(\texttt{kind},\texttt{payload}) \leftarrow \textsc{ParsePlanOrFinal}(\texttt{resp})$\;
  \If{$\texttt{kind}=\texttt{final}$}{
    $\texttt{final} \leftarrow \textsc{ValidateOrStringify}(\texttt{payload})$; merge $A$ and (optionally) $C$ into \texttt{final}\;
    \If{\textsc{RequestRequiresTools}(\texttt{user\_message}) and $T$ is empty and $b>0$ and $r_{\textsf{force}}<2$}{
      $r_{\textsf{force}}\leftarrow r_{\textsf{force}}+1$; append an instruction forcing a \texttt{tool\_call} response to $M$; \textbf{continue}\;
    }
    \KwRet{\texttt{final}, $T$}\;
  }
  \If{$\texttt{kind}\neq\texttt{tool\_call}$ or $b=0$}{Append to $M$ an instruction to return \texttt{final}; \textbf{continue}}\;

  $R \leftarrow [\ ]$ \tcp*[r]{compact results for planner}\;
  \ForEach{requested tool call $(\tau,\alpha)$ \textbf{while} $b>0$}{
    \If{$\tau=\texttt{cmap.colocalize}$}{$\alpha \leftarrow \textsc{SanitizeColocalizeArgs}(\texttt{user\_message},\alpha)$}\;
    $k \leftarrow (\tau,\textsc{CanonicalArgs}(\alpha))$\;
    \uIf{$k \in \mathcal{H}$}{
      $\texttt{compact} \leftarrow \mathcal{H}[k]$; record \texttt{cached} in $T$\;
    }
    \Else{
      $\texttt{result} \leftarrow \textsc{RunTool}(\tau,\alpha,\texttt{exec\_context})$\;
      normalize artifacts into $A$; append any emitted SDK code into $C$\;
      $\texttt{compact} \leftarrow \textsc{SummarizeForPlanner}(\tau,\texttt{result})$; set $\mathcal{H}[k]\leftarrow\texttt{compact}$\;
      record status/errors in $T$; $b\leftarrow b-1$\;
    }
    Append \texttt{compact} to $R$\;
  }
  Append a compact \texttt{TOOL\_RESULTS} observation containing $R$ to $M$\;
}
\end{algorithm}

\subsection{Validation and failure modes}
Inputs to computational tools are validated (e.g., numeric bounds, required parameters, tolerance ranges), and intermediate failures can be surfaced to the user with actionable messages.
This design supports iterative refinement, which is characteristic of real-world analysis and is emphasized by recent evaluations of data analysis agents \citep{Zhang2024DSEval,Li2025IDABench}.
